\section*{Bab \ref{ch:b2:plant-flowering} --- Perkembangan Reproduktif dan Kendali Pembungaan}

\begin{solution}{exo:b2:plant-flowering:1}
\omterm{def:b2:plant-meristems:meristem}{Meristem} vegetatif: membuat daun dengan tunas ketiak, tanpa batas (tak
tentu). \omterm{def:b2:plant-flowering:transition}{Meristem perbungaan}: membuat braktea dan \omterm{def:b2:plant-meristems:meristem}{meristem} \omterm{def:b2:angiosperm-reproduction:flower}{bunga} pada
ketiaknya, tak tentu pada banyak spesies. \omterm{def:b2:plant-meristems:meristem}{Meristem} \omterm{def:b2:angiosperm-reproduction:flower}{bunga}: membuat
keempat karangannya lalu berhenti --- tertentu.
\end{solution}

\begin{solution}{exo:b2:plant-flowering:2}
\omterm{prop:b2:plant-flowering:photoperiod}{Tumbuhan hari pendek} berbunga ketika harinya lebih pendek daripada
sebuah panjang kritis (krisan, kedelai, padi); \omterm{prop:b2:plant-flowering:photoperiod}{tumbuhan hari panjang}
ketika harinya lebih panjang (bayam, gandum, \emph{Arabidopsis});
tumbuhan netral hari berbunga begitu mencapai sebuah ukuran (tomat,
jagung). \omterm{prop:b2:plant-flowering:photoperiod}{Tumbuhan hari pendek} mengukur panjang malamnya yang tak
terputus.
\end{solution}

\begin{solution}{exo:b2:plant-flowering:3}
Tipe liarnya: kelopak, mahkota, \omterm{def:b2:angiosperm-reproduction:flower}{benang sari}, \omterm{def:b2:angiosperm-reproduction:flower}{daun buah}. Tanpa A: \omterm{def:b2:angiosperm-reproduction:flower}{daun
buah}, \omterm{def:b2:angiosperm-reproduction:flower}{benang sari}, \omterm{def:b2:angiosperm-reproduction:flower}{benang sari}, \omterm{def:b2:angiosperm-reproduction:flower}{daun buah}. Tanpa B: kelopak, kelopak,
\omterm{def:b2:angiosperm-reproduction:flower}{daun buah}, \omterm{def:b2:angiosperm-reproduction:flower}{daun buah}. Tanpa C: kelopak, mahkota, mahkota, kelopak, dan
\omterm{def:b2:angiosperm-reproduction:flower}{bunganya} berulang di dalamnya.
\end{solution}

\begin{solution}{exo:b2:plant-flowering:4}
Berpekan-pekan dingin dan lembap (musim dinginnya telah lewat); cahaya
(\omterm{def:b2:angiosperm-reproduction:seed}{bijinya} terletak dekat permukaannya); suhu yang berganti-ganti (tanah
terbuka, tanpa tajuk); pelucutan oleh hujan (hujannya telah datang);
penggoresan (perjalanan melewati sebuah usus atau pengikisan di dalam
tanah); asap dan bahang (sebuah kebakaran telah membersihkan tanahnya).
\end{solution}

\begin{solution}{exo:b2:plant-flowering:5}
16/8: malam 8 jam, di bawah 10 jam kritisnya: tidak. 12/12: ya.
12 terang, 6 gelap, sekilas cahaya, 6 gelap: dua malam masing-masing 6
jam: tidak. 8 terang, 4 gelap, 4 terang, 8 gelap: malam terpanjangnya 8
jam: tidak.
\end{solution}

\begin{solution}{exo:b2:plant-flowering:6}
\omterm{def:b2:plant-flowering:florigen}{Florigen} adalah sebuah isyarat yang dapat dipindahkan lewat cangkokan,
yang dibuat di dalam daunnya dan sama pada spesies hari pendek maupun
hari panjang --- yang berbeda adalah kapan daunnya membuatnya, bukan
apa yang dibuatnya. Dengan floemnya diputus, cangkokannya tidak
mengimbas pembungaannya: isyaratnya mengembara di dalam floemnya.
\end{solution}

\begin{solution}{exo:b2:plant-flowering:7}
$n = \ln(1/0.15)/0.05 = 1.90/0.05 = 38$ hari dingin. Gilirannya
berjumlah $47$: tumbuhannya tervernalisasi pada hari ke-11 giliran
ketiganya. Dengan $\alpha = 0.02$ ia akan membutuhkan 95 hari dan tidak
akan tervernalisasi.
\end{solution}

\begin{solution}{exo:b2:plant-flowering:8}
Gen penekannya dibungkam tanda kromatin yang disalin pada tiap
mitosisnya, sehingga tiap sel pucuknya mengingat musim dinginnya; di
dalam garis nutfah dan lembaga dininya tandanya dihapuskan dan gennya
diungkapkan lagi. Berguna karena sebutir \omterm{def:b2:angiosperm-reproduction:seed}{biji} yang dilepaskan pada
musim panas tidak boleh mewarisi izin tetuanya untuk berbunga: ia harus
menunggu musim dinginnya sendiri.
\end{solution}

\begin{solution}{exo:b2:plant-flowering:9}
M: berkecambah (fitokromnya diubah menjadi bentuk penyerap merah-jauh,
yaitu Pfr, bentuk yang aktif). M lalu MJ: tidak (diubah kembali menjadi
Pr). M, MJ, M: berkecambah. M, MJ, M, MJ: tidak. Hanya kilasan
terakhirnya yang berarti, karena tiap pengubahannya lengkap dan \omterm{def:b2:angiosperm-reproduction:seed}{bijinya}
membaca keadaan akhir pigmennya.
\end{solution}

\begin{solution}{exo:b2:plant-flowering:10}
$p = 0.7$: $r = 1.28$, $1.42$, $1.40$, $-\infty$ bagi $g = 0.4, 0.6, 0.8,
1$: terbaik di dekat $g = 0.7$. $p = 0.95$: $1.99$, $2.33$, $2.55$,
$-\infty$: optimumnya bergerak menuju $g = 1$ (kira-kira $0.948$),
hanya sebuah cadangan simbolis. Gurun, tempat tahun yang baik jarang,
memihak \omterm{def:b2:angiosperm-reproduction:seed}{dormansi} yang berat dan lumbung \omterm{def:b2:angiosperm-reproduction:seed}{biji} yang besar; padang rumput,
tempat sebagian besar tahunnya baik, lebih memihak pengecambahan
hampir segalanya.
\end{solution}

\begin{solution}{exo:b2:plant-flowering:11}
B pada semua karangannya: mahkota, mahkota, \omterm{def:b2:angiosperm-reproduction:flower}{benang sari}, \omterm{def:b2:angiosperm-reproduction:flower}{benang sari}.
C pada semua karangannya (C menekan A): \omterm{def:b2:angiosperm-reproduction:flower}{daun buah}, \omterm{def:b2:angiosperm-reproduction:flower}{benang sari}, \omterm{def:b2:angiosperm-reproduction:flower}{benang
sari}, \omterm{def:b2:angiosperm-reproduction:flower}{daun buah}. Tanpa A dan tanpa B: C saja di mana-mana --- \omterm{def:b2:angiosperm-reproduction:flower}{daun buah}
pada keempat karangannya. C juga mengakhiri \omterm{def:b2:plant-meristems:meristem}{meristemnya}; tanpanya
pusatnya tetap tak tentu lalu membuat \omterm{def:b2:angiosperm-reproduction:flower}{bunga} di dalam \omterm{def:b2:angiosperm-reproduction:flower}{bunga}.
\end{solution}

\begin{solution}{exo:b2:plant-flowering:12}
Panjang harinya diintegralkan sepanjang malamnya lewat keberimpitan
cahayanya dengan jamnya, dan pengimbasannya menuntut beberapa malam
semacam itu; dinginnya diintegralkan sebagai tanda kromatin yang
menimbun; lumbung \omterm{def:b2:angiosperm-reproduction:seed}{bijinya} mengintegralkan sepanjang tahun dengan
menyebarkan perkecambahannya. Sebuah pengintegral mengabaikan satu hari
yang menyimpang, sebuah pencairan Januari, atau satu tahun yang buruk;
sebuah ambang pada nilai seketikanya akan membungakan tumbuhannya pada
sebuah pekan hangat bulan Februari, atau mengecambahkan tiap \omterm{def:b2:angiosperm-reproduction:seed}{bijinya}
pada hujan yang pertama.
\end{solution}

\begin{solution}{pb:b2:plant-flowering:1}
\textbf{1.} Malam 8 jam lebih pendek daripada 9,5 jam kritisnya: tidak
berbunga. Untuk menjual pada bulan Agustus, penanamnya harus
memanjangkan malamnya secara buatan dengan kain hitam.
\textbf{2.} Malam 12 jam. Membuka pada 15 Agustus berarti
pengimbasannya selesai 8 pekan sebelumnya, yaitu sekitar 20 Juni;
kedua belas malam pengimbasnya harus mulai sekitar 8 Juni.
\textbf{3.} Sinarilah panennya selama beberapa menit di tengah tiap
malamnya: malam 14,5 jam menjadi dua malam masing-masing kira-kira 7
jam, keduanya di bawah panjang kritisnya. Sekilas cahaya sudah
mencukupi karena fitokromnya berubah dalam hitungan menit dan
tumbuhannya hanya mengukur panjang kegelapan yang tak terputus.
\textbf{4.} Sekilas cahaya pada pukul 21.00 menyisakan sebuah masa
gelap dari pukul 21.00 sampai 07.30, yaitu 10,5 jam --- lebih panjang
daripada 9,5 jam: ruas itu mengimbas. Selanya harus jatuh sedemikian
sehingga tidak ada ruas yang melampaui panjang kritisnya.
\textbf{5.} Satu malam yang terlewat memutus rangkaian dua belas malam
pengimbas yang berturut-turut: cacahnya mulai lagi dan pembungaannya
tertunda kira-kira sebanyak hari yang telah tertimbun. Sebatang
kokelbur membutuhkan satu malam pengimbas saja dan sudah akan
terpancang.
\textbf{6.} Fitokrom tidak menyerap cahaya hijau, sehingga tumbuhannya
tidak melihat kilasannya; merah mengubahnya menjadi bentuk aktifnya,
yang terbaca sebagai siang; merah-jauh segera sesudahnya mengubahnya
kembali, sehingga malamnya tidak terputus.
\textbf{7.} \omterm{prop:b2:plant-flowering:photoperiod}{Tumbuhan hari panjangnya} berbunga: \omterm{def:b2:plant-flowering:florigen}{florigen} menyeberangi
sambungan cangkokannya di dalam floemnya. Hal itu mengenali sebuah
isyarat yang berpindah, dapat dipindahkan lewat cangkokan, dan sama
bagi kedua golongan panjang harinya.
\textbf{8.} $n = \ln 10/0.06 = 38.4$: yaitu 39 hari dingin.
\textbf{9.} 10 hari pada bulan November, 35 hari sampai akhir Januari,
dan hari dingin ke-39 jatuh pada hari dingin keempat bulan Februari:
tervernalisasi pada awal Februari.
\textbf{10.} Enam hari: $f = e^{-0.36} = 0.70 > 0.10$: ia tetap
vegetatif sepanjang musim panas dan tidak memberi bulir --- gandum
musim dingin harus ditanam pada musim gugur.
\textbf{11.} Tanpa penekannya, ia berbunga tanpa dingin: ia dapat
ditanam pada musim semi lalu dipanen pada musim panas yang sama ---
yaitu gandum musim semi.
\textbf{12.} Tanda pembungaman pada kromatin penekannya disalin pada
tiap mitosisnya, sehingga pucuk apikalnya mengingat sepanjang musim
semi; pada \omterm{def:b2:meiosis-heredity:meiosis}{meiosis} dan lembaganya, tandanya dihapuskan, dan tiap
generasinya harus mencacah musim dinginnya sendiri.
\textbf{13.} Ingatannya harus duduk di dalam sel yang akan membuat
\omterm{def:b2:angiosperm-reproduction:flower}{bunganya} dan yang bertahan sepanjang musim dingin --- yaitu \omterm{def:b2:plant-meristems:meristem}{meristemnya}
--- dan ia adalah sebuah keadaan kromatinnya, bukan sebuah isyarat yang
berpindah; daun yang mengukur panjang harinya rontok, dan mengekspor
sebuah protein sebagai gantinya.
\textbf{14.} Golongan C: kelopak, mahkota, mahkota, kelopak (A meluas ke
karangan keempatnya dan \omterm{def:b2:angiosperm-reproduction:flower}{bunganya} berulang di dalamnya).
\textbf{15.} Tanpa \omterm{def:b2:angiosperm-reproduction:flower}{benang sari}, maka tanpa \omterm{def:b2:plant-life-cycles:heterospory}{serbuk sari}; beberapa \omterm{def:b2:angiosperm-reproduction:flower}{daun
buah} terbentuk di tempat C-nya lemah aktif, sehingga sesekali \omterm{def:b2:angiosperm-reproduction:seed}{bijinya}
terjadi.
\textbf{16.} Golongan A: \omterm{def:b2:angiosperm-reproduction:flower}{daun buah}, \omterm{def:b2:angiosperm-reproduction:flower}{benang sari}, \omterm{def:b2:angiosperm-reproduction:flower}{benang sari}, \omterm{def:b2:angiosperm-reproduction:flower}{daun buah}.
\textbf{17.} C menekan A: karangan 1 menjadi \omterm{def:b2:angiosperm-reproduction:flower}{daun buah}, karangan 2 (B
dan C) menjadi \omterm{def:b2:angiosperm-reproduction:flower}{benang sari} --- \omterm{def:b2:angiosperm-reproduction:flower}{daun buah}, \omterm{def:b2:angiosperm-reproduction:flower}{benang sari}, \omterm{def:b2:angiosperm-reproduction:flower}{benang sari},
\omterm{def:b2:angiosperm-reproduction:flower}{daun buah}, yaitu \omterm{def:b2:angiosperm-reproduction:flower}{bunga} mutan A.
\textbf{18.} Gen ABC-nya adalah satu keluarga \omterm{prop:b2:cell-differentiation:transcriptional}{faktor transkripsi} yang
diwarisi dari moyang bersama semua tumbuhan berbunga; modulnya sudah
ada sebelum mawar, snapdragon, dan \emph{Arabidopsis} berpisah, dan
masing-masing mempertahankannya.
\textbf{19.} B tanpa A maupun C tidak menetapkan organ \omterm{def:b2:angiosperm-reproduction:flower}{bunga} apa pun:
sebuah organ mirip daun atau mirip braktea.
\textbf{20.} $p = 0.8$: $r = 1.44$ ($g = 0.5$), $1.59$ ($0.7$), $1.55$
($0.9$).
\textbf{21.} Dari ketiganya, $g = 0.7$ (optimum sejatinya terletak di
dekat $0.8$); pada $g = 1$ faktor tahun buruknya nol dan garis
keturunannya dimusnahkan tahun buruk yang pertama.
\textbf{22.} $p = 0.5$: $0.51$, $0.41$, $-0.03$: $g = 0.5$ yang
terbaik, dan yang lebih rendah lagi akan lebih baik --- \omterm{def:b2:angiosperm-reproduction:seed}{dormansi} yang
lebih banyak di gurun.
\textbf{23.} \omterm{def:b2:angiosperm-reproduction:seed}{Biji} penanamnya disimpan, bukan dikubur, dan bertahan
hidup dengan kepastian; ia tidak perlu berlindung di dalam ladangnya,
tetapi harus menyimpan sebuah cadangan di dalam lumbungnya untuk
ditanam ulang sesudah sebuah tahun yang gagal --- cadangan yang sama,
yang dipegangnya alih-alih dipegang tanahnya.
\textbf{24.} Sebutir \omterm{def:b2:angiosperm-reproduction:seed}{biji} yang kecil mengusung cadangan bagi sebuah
pucuk sepanjang satu atau dua sentimeter; kalau berkecambah di dalam
gelap di bawah tanah, ia akan mati sebelum mencapai cahaya, sehingga ia
menunggu cahaya --- yang hanya mencapainya dekat permukaannya.
Pembajakan membawa \omterm{def:b2:angiosperm-reproduction:seed}{biji} yang terkubur ke atas lalu memicu ledakan
gulma.
\textbf{25.} Malam kritisnya 9,5 jam, penutupannya dari sekitar 8 Juni
agar membuka pada 15 Agustus; 39 hari dingin; nilai terbaiknya
$g \approx 0.8$ bagi $p = 0.8$ dan kira-kira $0.5$ bagi $p = 0.5$.
\end{solution}
